% RESEARCH_DRAFT_SHA256: 811c68a28cb78c2e0f8d2af436525d7d4a4d21d5d9db052890ad7afa6a55b34b
% Layout and source-reference wrappers expanded for the web renderer.
\documentclass[11pt]{article}
\usepackage[margin=25mm]{geometry}
\usepackage{fontspec,unicode-math}
\setmainfont{Latin Modern Roman}
\setmathfont{Latin Modern Math}
\usepackage{amsmath,amssymb,mathtools,xurl,hyperref}
\hypersetup{colorlinks=true,urlcolor=blue,linkcolor=blue}
\setcounter{secnumdepth}{0}
\setlength{\parindent}{0pt}
\setlength{\parskip}{5pt}
\setlength{\emergencystretch}{3em}
\newcommand{\sref}[2]{\hyperlink{source:#1:#2}{[S#2]}}
\newcommand{\cataloguescope}[1]{\paragraph{Source record.}#1}
\newcommand{\floor}[1]{\left\lfloor#1\right\rfloor}
\begin{document}
\section*{Research-attempt provenance}
This standalone research attempt retains the complete original definition body below,
including its historical source and preparation notes. Those notes describe the
earlier catalogue preparation; the new attempt and its verification are identified
separately after that body. No original ID, catalogue status, or source statement is changed.

\section{Predicting SAT-Solver Hardness from Formula Structure}\label{statement}

\subsection{Definitions and mathematical statement}

\paragraph{Shared notation.}Standard number systems are $\mathbb N,\mathbb Z,\mathbb Q,\mathbb R,\mathbb C$. Specialized notation and index ranges follow the cited source and the definitions below.


A formula-space measure bounds the maximum number of clauses or formulas retained simultaneously during a derivation. It differs from proof length and depends on the chosen proof system.
\paragraph{Record scope and qualifications.}
The catalogue title is retained as metadata. Named certificates, contracts, bounds or frameworks in that title are not asserted as results of the original source extract.
Hardness may mean worst-case complexity or performance of one solver; the supplied question does not choose a solver, distribution, or quantitative hardness criterion.

Source-record limitation: The explanation of solver success and formula hardness is a broad agenda without a solver model, hardness criterion, or stated proposition.
\paragraph{Statement recorded in the supplied source.}
What explains the success of SAT solvers, how can one determine whether a particular propositional formula is hard or tractable, and can formula space complexity serve as a useful measure of hardness? \href{https://doi.org/10.4171/owr/2011/52}{[S2]}
\paragraph{Source status.}
The supplied catalogue labels this source record partially solved; that metadata was checked against the supplied snapshot on 2026-10-11. The label is preserved for provenance and is not an independent verification of present open status or mathematical correctness. The catalogue literature assessment is dated 2026-08-21.
\subsection{Short English statement}
Explain satisfiability-solver difficulty using formula structure and assess the role of proof space.
\subsection{Sources}
\begin{itemize}
\item[S1] Ulam AI and the UnsolvedMath Contributors, UnsolvedMath: A Curated Collection of Open Mathematics Problems, catalogue entry. Catalogue curation and metadata: CC BY 4.0. The source extract is retained; definitions, English summaries and qualifications are editorial additions.
\url{https://huggingface.co/datasets/ulamai/UnsolvedMath}.
\item[S2] Mathematical Logic: Proof Theory, Constructive Mathematics (2012), Oberwolfach Reports, DOI 10.4171/owr/2011/52, pp. 2963--3002
\url{https://doi.org/10.4171/owr/2011/52}.
\item[S3] Ben-Sasson--Nordström
\url{https://arxiv.org/abs/1008.1789}.
\item Source-specific Oberwolfach Reports licences, including CC BY-SA 4.0 where applicable, are retained; see THIRD\_PARTY\_NOTICES.md.
\end{itemize}
\paragraph{Preparation and review.}
The original source and UnsolvedMath contributors retain attribution for the source extract. Codex (GPT-6; exact serving revision not exposed) assisted with definitions, English summaries, scope notes and formatting. No new human mathematical authorship or independently verified solution is asserted. Catalogue identity checked against the complete supplied registry and current repository. Distinct upstream records, including repeated questions and broad or damaged extracts, are retained for complete coverage of the source CS category. A repeated record is not a new independent conjecture. Complete mathematical-equivalence and current-literature review has not been performed. Preparation and final review follow contribution prompts C and D, with source gaps disclosed. Structural and rendering checks do not establish a new theorem or open status.
\label{end:statement}

\section{Literature and context check}
The primary Nordstrom report, printed page 2991, proposes proof space as a useful hardness measure and relates time-memory tradeoffs to pebbling. This attempt does not import those lower-bound constructions. It fixes resolution clause space for an explicit family and separately fixes a simple search algorithm to show why a solver model matters. These are limited tests of a broad explanatory agenda.
\section{Research attempt: A clause-space calculation and solver-dependent search examples}

\paragraph{Constant clause space with unbounded proof length}
For $n\geq1$ use the unsatisfiable chain
\[
 F_n=\{x_1,\ \neg x_1\vee x_2,\ldots,\neg x_{n-1}\vee x_n,\ \neg x_n\}.
\]
In resolution, a configuration stores clauses; an inference adds a resolvent of two currently stored clauses, and clauses may be erased. Count input-clause downloads and inferences in proof length, and the maximum number of stored clauses in clause space. Starting from the unit $x_1$, download $\neg x_i\vee x_{i+1}$, resolve to $x_{i+1}$, and erase the old unit and binary clause. At most three clauses are stored. Finally download $\neg x_n$ and derive the empty clause, again within space three. This gives a linear-length, constant-space refutation.

Every input clause is essential to unsatisfiability. Removing $x_1$ permits the all-zero assignment; removing $\neg x_n$ permits all ones. Removing the implication $\neg x_i\vee x_{i+1}$ permits $x_1=\cdots=x_i=1$ and $x_{i+1}=\cdots=x_n=0$. Thus any resolution refutation must download all $n+1$ input clauses: otherwise it would refute a satisfiable subset. Resolution soundness used here follows directly because any assignment satisfying both $A\vee x$ and $B\vee\neg x$ satisfies $A\vee B$; induction propagates satisfaction through all inferences. Consequently proof length is at least $n+1$ despite the uniform space-three upper bound.

This is not a difficult SAT family; unit propagation solves it efficiently. Its purpose is to show that clause space alone, without formula size or the proof system, cannot upper-bound even elementary reading/proof length by a constant depending only on that space.

\paragraph{A fixed search order can dominate superficial structure}
Consider exhaustive search that tests assignments in lexicographic order with zero before one and stops at the first satisfying assignment. The formulas
\[
 A_n=\bigwedge_{i=1}^n\neg x_i,\qquad B_n=\bigwedge_{i=1}^n x_i
\]
have the same numbers of variables, clauses and literal occurrences and the same clause width. They differ only by globally complementing all variables. The specified search checks one assignment for $A_n$ and all $2^n$ assignments for $B_n$, because each has a unique satisfying assignment at the corresponding end of the order. A unit-propagation solver, however, fixes every variable immediately in either formula and uses only linear work in their explicit size. This is an exact comparison for the two stated algorithms, not a claim about modern solver heuristics.

\paragraph{The missing general measure}
For unsatisfiable inputs, proof-space lower bounds refer to a chosen proof system and constrain algorithms whose executions yield proofs in that system. Satisfiable inputs instead need a witness search criterion; resolution refutation space is not defined as a finite refutation resource for them. The examples above also show that input size and algorithmic invariances cannot be discarded. To turn the report's proposal into a predictive theorem one must specify the solver, representation, input distribution or worst-case quantifiers, and the relation between memory space and elapsed search. No such universal relation is proved by these elementary families.

\section{Verification and adversarial audit}
The clause-space convention counts the newly inferred clause before old premises are erased, giving space three rather than an unjustified two. Clause necessity is checked by explicit satisfying assignments after each deletion. The exhaustive-search count is for a precisely stated ordering and is contrasted with a proved unit-propagation behavior on the same formulas.
\paragraph{Audit corrections.} Specified the proof-space convention, counted clause downloads, and kept satisfiable search complexity separate from unsatisfiable refutation resources.
\paragraph{External computation.} None; symbolic arguments and boundary checks only.
\section{Scope and remaining gaps}
The broad question of explaining SAT-solver success and predicting hardness remains unresolved. These exact resource calculations delimit what a formula-space measure can claim without additional size and solver information; they do not rederive the source's pebbling-based tradeoff results.
\paragraph{Integration review.} Codex (GPT-6; exact serving revision not exposed) reviewed the written arguments and expanded layout and source-reference wrappers. Repository metadata and the production web renderer were checked using Python 3.13 and JavaScript with MathJax 3.2.2. This model-assisted review is not independent human proof verification; no PDF compilation or formal proof checking is asserted.

\section{Completion Estimate}
\noindent\textbf{COMPLETION: 0\%}
The calculations below the preserved source statement investigate a concrete route,
but establish no substantiated new progress toward the entire original target.
Known special cases and failed sufficient conditions are not credited as a solution.
\section{References}
The original source references and catalogue attribution are preserved in the
Sources section above. The following URLs were actually consulted for this attempt;
retention of an additional original reference does not claim fresh consultation of it.
\begin{itemize}
\item \url{https://ems.press/content/serial-article-files/46368} (consulted 2026-10-11). 
\end{itemize}
\section{Model and contributor disclosure}
The parent Codex GPT-6 model has reviewed the written mathematical arguments at their declared scopes. This is model-assisted review, not independent human verification.
OpenAI gpt-6-astra, requested reasoning configuration ultra, generated this research
attempt on 2026-10-11. The exact serving revision is not exposed. Codex GPT-6
performed orchestration and repository-format review; its exact serving revision
is also not exposed. The mathematical route, boundary checks and final adversarial
self-audit followed contribution prompts A then B in
\texttt{docs/contribution\_prompts.md}, after reading \texttt{CONTRIBUTING.md}.
Python 3.14.2 assembled the standalone documents; any mathematical computation is
specified in Verification. No theorem prover was used. No human mathematical
authorship, substantial human proof correction, or independent human proof checking
is asserted. Model-assisted review, source retrieval, and structural checks do not
establish correctness or novelty. Full private reasoning transcripts are unavailable
as reproducibility artifacts; the complete arguments, inputs and outcomes actually
used are stated here or in the companion batch audit.
\section{Final status}
Unresolved as a new research contribution. This is a concrete route investigation,
not a claim of new complete resolution or a current-status certification. A known
answer reproduced for context does not become a new contribution. The catalogue's
source-status label is preserved; this attempt status does not assert that a
source-labelled solved question is open.
\end{document}
